\documentclass[11pt]{article}

\usepackage[margin=1in]{geometry}
\usepackage[T1]{fontenc}
\usepackage[utf8]{inputenc}
\usepackage{lmodern}
\usepackage{enumitem}
\usepackage{booktabs}
\usepackage{xcolor}
\usepackage[most]{tcolorbox}
\usepackage[hidelinks]{hyperref}

\setlength{\parindent}{0pt}
\setlength{\parskip}{6pt}
\setlist[itemize]{leftmargin=1.6em,itemsep=0.25em,topsep=0.2em}

\definecolor{commentgray}{RGB}{235,235,235}
\definecolor{changeblue}{RGB}{0,0,180}

\newtcolorbox{reviewercomment}{
  enhanced,
  breakable,
  colback=commentgray,
  colframe=commentgray,
  boxrule=0pt,
  arc=0pt,
  left=5pt,
  right=5pt,
  top=5pt,
  bottom=5pt,
  before skip=10pt,
  after skip=10pt,
  fontupper=\itshape
}

\newcommand{\response}{\noindent\textbf{Response:} }
\newcommand{\manuscriptchange}[1]{%
  \par\noindent{\color{changeblue}\textbf{Manuscript change:} #1}\par
}

\begin{document}

\begin{center}
{\LARGE Response to Reviewers}\\[6pt]
{\large Journal of Visual Communication and Image Representation --- Manuscript JVCI-26-1814}\\[6pt]
{\large Net-Aware Learned Block Mapping for Reversible Data Hiding in Binary Images}\\[18pt]
August 20, 2026
\end{center}

We thank the Senior Area Editor and the reviewers for their detailed assessment. The revision substantially refocuses the manuscript around a single defensible technical contribution: \emph{serialization-aware net-capacity-driven activation of reversible PEAK/ZERO mappings}. We now state explicitly that net-capacity accounting, binary RDH, PEAK/ZERO block mapping, side information, and distortion-aware ranking are established concepts. The algorithmic distinction is that the exact serialized description cost required for deterministic recovery is used inside the active mapping-selection decision, rather than subtracted only after gross capacity has been maximized.

The article is written as a self-contained work. We corrected the bibliography, narrowed the claims, added the requested ablations, and strengthened external validation. We also revised the scope language so that three conceptually different issues are no longer conflated: \emph{conditions of validity} of the exact-recovery contract, \emph{operating characteristics} measured within that contract, and \emph{true limitations} of the evidence or intended use.

\textbf{Overview of major changes.}

\begin{itemize}
  \item \textbf{Novelty refocused:} the core claim is serialization-aware PEAK/ZERO activation; net accounting itself is treated as established.
  \item \textbf{Wire model clarified:} the revision adds encoder/decoder contracts and an executable auxiliary-stream table for ABM and baselines.
  \item \textbf{Ablations added:} A0/A1/A2 activation, non-learning rankers, and CNN architecture diagnostics now isolate the useful contribution.
  \item \textbf{Evaluation strengthened:} confidence intervals, runtime decomposition, sensitivity analysis, residual visuals, channel stress, and a frozen external FUNSD validation are now reported.
  \item \textbf{External document validation added:} the committed CNN-ranked codec is evaluated without retraining or tuning on all 199 real scanned FUNSD documents, under fixed-128 and Otsu binarization and at 64, 128, 256, and 512 bits. Every available run recovers both message and cover exactly; A2 is never worse than A1 in paired net payload and saves approximately 239--241 auxiliary bits per paired image.
  \item \textbf{Scope classified:} matched stego/auxiliary data and lossless transport are conditions of validity; metadata-dominated low-payload behavior, availability, DRD, and payload-dependent detectability are operating characteristics; the remaining true limitations include high-payload detectability for covert use, protocol dependence of baseline serialization, and external breadth beyond one independent scanned-document corpus.
\end{itemize}

\section*{Response to the Senior Area Editor}

\begin{reviewercomment}
The authors are suggested to revise the manuscript according to the reviewer comments point by point carefully, especially clarifying the contributions and technique details, justifying the scheme designing and strengthening the experiment analysis.
\end{reviewercomment}

\response We agree. The revised manuscript presents the proposed scheme through one focused mechanism: serialization-aware activation of reversible PEAK/ZERO mappings. We added the missing technical contract for the auxiliary stream, the exact wire representation, the central activation ablation, non-learning ranker baselines, CNN architecture diagnostics, sensitivity analysis, runtime decomposition, visual residuals, channel stress tests, and a frozen external validation on all 199 FUNSD scanned documents without retraining or tuning.

\manuscriptchange{The main revisions appear in the Abstract, the Introduction through Conclusion, the wire-model table, the activation-ablation table, the ranker-ablation table, the CNN-architecture table, the new FUNSD external-validation table, the module-runtime table, the channel-stress table, and the multi-payload, protocol-overview, steganalysis, and visual-comparison figures.}

\section*{Response to Reviewer \#1}

\begin{reviewercomment}
1. The core mechanism is conventional $3\times3$ binary PEAK/ZERO block mapping; candidate ranking and auxiliary-bit accounting do not by themselves constitute a fundamentally new RDH mechanism.
\end{reviewercomment}

\response We agree with the premise and revised the novelty claim accordingly. The manuscript now states that the reversible carrier family is established and deterministic. The contribution is the selection rule that activates only those reversible mappings whose exact serialized recovery cost still leaves usable mapping capacity.

\manuscriptchange{The Abstract, Introduction, Related Work, and Discussion now use the phrase ``serialization-aware net-capacity-driven activation'' and treat block mapping, RDH, and net accounting as established foundations.}

\begin{reviewercomment}
2. Auxiliary information must be counted in RDH. Subtracting synchronization overhead should be treated as necessary evaluation practice, not as the main contribution.
\end{reviewercomment}

\response We agree. The manuscript now treats net-capacity accounting itself as established. The distinction is that the exact serialized cost of the reversible mapping participates in deciding which PEAK/ZERO tables remain active.

\manuscriptchange{The Introduction and proposed-method section now separate established net-capacity accounting from the proposed serialization-aware mapping-selection mechanism.}

\begin{reviewercomment}
3. The CNN ranking gives a limited DRD reduction and does not justify making learned block mapping the main emphasis.
\end{reviewercomment}

\response We agree that the CNN serves as a supporting component. The revised manuscript treats it as a learned distortion ranker supporting the reversible mechanism. The central novelty claim is moved to serialization-aware activation. The CNN results are retained as a secondary ablation.

\manuscriptchange{The Abstract, contribution list, and Conclusion now identify the CNN as a secondary learned candidate-cost ranker rather than the core reversible mechanism.}

\begin{reviewercomment}
4. Compare the CNN against simpler non-learning rankers such as direct DRD-cost, transition-preserving, or connectivity-aware costs.
\end{reviewercomment}

\response We agree. We added a ranker ablation using identical distance-one candidates and a common payload/serialization protocol. The CNN remains close to direct DRD-cost ranking and improves over Hamming-index and simple transition/connectivity heuristics, while the central net-payload gain comes from serialization-aware activation.

\manuscriptchange{The ``Candidate Ranking and CNN Architecture'' subsection and the ranker-ablation table report the Hamming, direct DRD, transition-preserving, connectivity-aware, and CNN rankers.}

\begin{reviewercomment}
5. The CNN only orders candidate ZERO patterns; it does not learn an end-to-end reversible process.
\end{reviewercomment}

\response We agree and now state this explicitly. The deployed codec remains a deterministic distance-one PEAK/ZERO reversible substitution system. The CNN supplies an ordering over admissible substitutions; invertibility is provided by the explicit mapping and wire representation.

\manuscriptchange{The proposed-method section now states that the CNN is a candidate ranker and that the deterministic mapping provides invertibility.}

\begin{reviewercomment}
6. Gaussian-process exploration and the Random Forest uniform-block layer are diagnostic, not deployed components.
\end{reviewercomment}

\response We agree. Those parts have been condensed and relabelled as diagnostic ablations. The deployed codec is defined by distance-one PEAK/ZERO substitution, serialization-aware net-capacity activation, enumerative auxiliary coding, and secondary CNN ranking as a deployed cost-ordering component.

\manuscriptchange{The Introduction and Section ``Learned Ranking, Ablations, and Wire Accounting'' now identify Gaussian-process and Random Forest analyses as diagnostic only.}

\begin{reviewercomment}
7. ABM has negative mean net payload at 64 and 128 bits, which limits practical robustness.
\end{reviewercomment}

\response We agree that the negative net values are important and must be reported, but the revised manuscript distinguishes this behavior from a failure of validity. At low requested payloads, the fixed and mapping-dependent synchronization cost exceeds the application payload; exact recovery still holds whenever the target is available. We therefore classify this as a \emph{metadata-dominated operating characteristic} of the tested payload regime rather than a limitation of the reversible contract. The new FUNSD experiment reinforces this distinction: exact recovery remains universal among available runs, while metadata overhead continues to dominate delivered net payload through 512 bits on that scanned-document corpus.

\manuscriptchange{Sections ``Frozen External Validation on FUNSD,'' ``Multi-Payload Common-Protocol Results,'' ``Positioning Across Practical Objectives,'' and ``Experimental Scope'' now report negative-net low-payload behavior as an operating characteristic and separate it from conditions of validity and true limitations.}

\begin{reviewercomment}
8. Logistic detector AUC values become high at 256 and 512 bits, limiting security relevance.
\end{reviewercomment}

\response We agree. Unlike metadata amortization, this is retained as a true limitation for covert or security-prioritized use under the tested detector. The revised text positions the validated high-payload region as capacity-oriented reversible annotation over a lossless channel and reports the AUC values directly. We also added a quantitative detectability-cause analysis using block-histogram Jensen--Shannon displacement and horizontal/vertical transition changes.

\manuscriptchange{The Abstract, Grouped Steganalysis subsection, Experimental Scope, and Conclusion now identify high-payload detectability as a genuine limitation on covert-use claims and analyze why it rises.}

\begin{reviewercomment}
9. The auxiliary stream is ambiguous: embedded in the stego image, transmitted through a side channel, or stored externally.
\end{reviewercomment}

\response We agree. The revised manuscript now defines the codec contract as $(I,M)\mapsto(I',A_{\rm wire})$ and $(I',A_{\rm wire})\mapsto(I,M)$. The auxiliary stream is serialized side information stored or transmitted alongside the marked image in the evaluated implementation. Preservation of a matched marked image and valid auxiliary stream over a lossless path is now stated explicitly as a \emph{condition of validity} for exact recovery, not as an empirical limitation of an otherwise matched execution.

\manuscriptchange{The ``Overview and Design Rationale'' subsection adds the encoder/decoder contract, the wire-model table gives the exact charged fields, and ``Experimental Scope'' now labels matched stego/auxiliary data and lossless transport as conditions of validity.}

\section*{Response to Reviewer \#2}

\begin{reviewercomment}
1. Clarify which components are essential and which are diagnostic.
\end{reviewercomment}

\response We agree. The revised contribution hierarchy is: essential mechanism, serialization-aware PEAK/ZERO activation; executable realization, enumerative wire representation; secondary distortion module, CNN ranking; diagnostics, Gaussian-process exploration and Random Forest uniform-block analysis; evaluation contribution, common executable protocol.

\manuscriptchange{The contribution list in the Introduction and the method overview now separate deployed components from diagnostics.}

\begin{reviewercomment}
2. The CNN improvement is limited; justify its necessity.
\end{reviewercomment}

\response We agree that the CNN should be justified as a useful supporting component. The new ranker ablation shows that CNN ranking is a distortion-ranking component, while the principal net-payload gain comes from serialization-aware activation. The proposed codec remains executable with non-learning rankers.

\manuscriptchange{The ranker-ablation and CNN-architecture tables now quantify CNN utility and architecture sensitivity.}

\begin{reviewercomment}
3. Baseline net-payload comparison may depend on authors' serialization choices.
\end{reviewercomment}

\response We agree and retain this as a true limitation of cross-method net-payload inference. The revised wire-model table states which parts are specified in the cited papers and which executable representations were introduced for the common protocol. We also added a fixed-header serialization-sensitivity analysis that credits baselines with 104 auxiliary bits while retaining payload-critical maps, states, locations, and payload length. The conclusion is limited to the tested executable representations and fixed-header variation.

\manuscriptchange{The wire-model table, the ``Serialization Sensitivity for Baselines'' subsection, ``Experimental Scope,'' and the Conclusion now state this protocol-dependence limitation directly.}

\begin{reviewercomment}
4. The method has a clear detectability problem at higher payloads.
\end{reviewercomment}

\response We agree. High-payload detectability is treated as a true limitation for covert use. The method is scoped as reversible annotation over lossless channels, and the Abstract and Conclusion position the high-payload settings as capacity-oriented rather than covert-security operating points under the tested detector.

\manuscriptchange{The Abstract, Section ``Grouped Steganalysis,'' Experimental Scope, and Conclusion were revised to narrow the security claim.}

\begin{reviewercomment}
5. The learning-assisted method is trained on eight document images and transferred to BOSSbase; more generalization detail is needed.
\end{reviewercomment}

\response We agree and have added a new frozen external validation specifically to address this concern. The committed CNN ranker is applied, without retraining, fine-tuning, threshold fitting, or other tuning on FUNSD, to all 199 real scanned documents in the FUNSD corpus. Each document is independently evaluated under fixed threshold 128 and Otsu binarization at 64, 128, 256, and 512 bits. Every available run recovers the payload exactly and restores the binary cover bit-for-bit. Availability ranges from 98.99\% to 100\% across the tested binarization/activation/payload conditions. On paired available images, A2 is never worse than A1 in delivered net payload and reduces auxiliary cost by a mean of 238.9--241.4 bits, with a median saving of 232 bits. This materially addresses the previous absence of a sizeable real scanned-document external validation and shows transfer of both exact reversibility and the proposed activation benefit without retraining.

We have deliberately not converted this result into a broader efficiency claim. Mean net payload remains negative through 512 bits on FUNSD, and the largest positive-net fraction is 4/198 (2.02\%). The remaining external-generalization limitation is therefore no longer the absence of any real scanned-document benchmark; it is the breadth of evidence beyond one independent scanned-form corpus, including other document collections, scanner pipelines, languages, and acquisition conditions.

\manuscriptchange{Sections ``Datasets, Splits, and Reproducibility,'' ``Frozen External Validation on FUNSD,'' ``Experimental Scope,'' and the Conclusion now report the frozen FUNSD protocol, exact recovery, A1/A2 paired results, and the narrower remaining generalization limitation.}

\section*{Response to Reviewer \#3}

\begin{reviewercomment}
1. Provide more contribution detail in the abstract, not only many results.
\end{reviewercomment}

\response We agree. The abstract was rewritten to foreground the mechanism: serialization-aware activation, exact wire representation, CNN as ranker, common protocol, and frozen real-document transfer. It also distinguishes the demonstrated transfer of exact reversibility/activation benefit from the metadata-dominated net-efficiency regime observed on FUNSD.

\manuscriptchange{The Abstract was replaced and aligned with the revised contribution hierarchy and external-validation evidence.}

\begin{reviewercomment}
2. Steganalysis is usually for steganography rather than RDH; reorganize or remove those sections.
\end{reviewercomment}

\response We agree that steganalysis should serve the RDH scope. We retained a shorter detector analysis because the reviewers specifically raised detectability as an operating concern, and we moved its interpretation to scope. High-payload detectability is now identified specifically as a true limitation on covert/security-prioritized use, rather than treating the detector as a deployed security module.

\manuscriptchange{The steganalysis material is presented as detectability/scope evaluation in Section ``Grouped Steganalysis'' and the Discussion.}

\begin{reviewercomment}
3--4. Add visual quality and visualization comparisons including cover, marked image, and residual maps for all methods.
\end{reviewercomment}

\response We agree. We added a visual comparison figure showing matched crop-level cover, stego/marked images, and residual maps for ABM, Dong, Huynh--Nguyen, and PPOCP under the same payload.

\manuscriptchange{The visual-comparison figure adds the requested cover/stego/residual visualization.}

\section*{Response to Reviewer \#4}

\begin{reviewercomment}
1. Provide computational cost analysis by module, not only overall runtime.
\end{reviewercomment}

\response We agree. The revised manuscript includes a module-level runtime decomposition for candidate generation/ranking, prefix optimization, embedding, serialization, extraction, and verification.

\manuscriptchange{The ABM module-runtime table reports the runtime decomposition.}

\begin{reviewercomment}
2. Discuss payload size, auxiliary length, and mapping parameters through sensitivity analysis.
\end{reviewercomment}

\response We agree. We added payload/auxiliary sensitivity analysis, including how net payload, DRD, histogram displacement, and transition changes evolve with payload fraction; changed-pixel statistics are provided in the reproducibility artifacts. We also added serialization sensitivity for baselines. The new FUNSD experiment provides an additional distribution-level sensitivity result: serialization-aware A2 reduces auxiliary cost systematically under both fixed-128 and Otsu binarization, while the fixed synchronization burden keeps delivered net payload metadata-dominated through 512 bits. We classify this as an operating characteristic rather than a validity failure.

\manuscriptchange{Sections ``Serialization Sensitivity for Baselines,'' ``Frozen External Validation on FUNSD,'' and ``Multi-Payload Common-Protocol Results'' now discuss parameter influence and the operating envelope.}

\begin{reviewercomment}
3. Define parameters such as $\mathrm{HT}_i=\mu_i+\alpha\sigma_i$ when first introduced.
\end{reviewercomment}

\response We agree. The adaptive threshold equation now defines $\mu_i$, $\sigma_i$, $\alpha$, and the resulting candidate set immediately where the notation appears. The adaptive-threshold parameter $\alpha$ belongs to the unconstrained diagnostic exploration; the deployed codec fixes the policy to Hamming distance one.

\manuscriptchange{The adaptive-threshold and adaptive-candidate-set equations now define the threshold and candidate set explicitly.}

\begin{reviewercomment}
4. Add more analysis of why methods differ in net capacity and distortion.
\end{reviewercomment}

\response We agree. The revised Discussion decomposes the observed differences into reversible mapping capacity, serialized synchronization cost, distortion ranking, and detectability footprint. The new central ablation isolates the net-payload effect of activation itself, and FUNSD shows that the same A2 synchronization saving persists under an external scanned-document distribution even when absolute net payload remains negative at the tested targets.

\manuscriptchange{The Discussion, activation-ablation table, and FUNSD external-validation subsection now provide this decomposition.}

\section*{Response to Reviewer \#5}

\begin{reviewercomment}
1. Justify the selected CNN architecture with ablations over depth, width/window, MAE, and inference or training cost.
\end{reviewercomment}

\response We agree. We added a compact architecture diagnostic comparing shallow, current, deeper, different-window, and different-width variants under the same 20-epoch protocol, optimizer, learning rate, batch size, loss, split, and seed policy. The table now includes validation MAE, Pearson correlation, held-out placement DRD, training time, and inference time per candidate. The wider two-layer 9-by-9 model gives the lowest MAE, while the deployed two-layer 9-by-9 configuration remains a lighter compromise with close MAE, lower inference cost, and comparable placement DRD.

\manuscriptchange{The CNN-architecture table now reports depth, window-size, and width diagnostics with inference timing.}

\begin{reviewercomment}
2. The Random Forest uniform-block layer is mathematically non-viable in net payload and distracts from the method.
\end{reviewercomment}

\response We agree. The revised manuscript removes it from the core method and keeps it only as a diagnostic example showing why gross-capacity additions can be invalidated by coordinate overhead.

\manuscriptchange{The Random Forest material is condensed and labelled as a diagnostic ablation in the proposed-method section.}

\begin{reviewercomment}
3. High-payload AUC is severe; either implement a security-aware ranker or establish a quantified payload threshold.
\end{reviewercomment}

\response We agree. We evaluated the payload-threshold option and the measured AUC values support a conservative scope decision: the manuscript designates no tested payload as a covert-security threshold. The 256- and 512-bit configurations, with AUC 0.900 and 0.957, are explicitly restricted to capacity-oriented reversible annotation. The 64- and 128-bit settings exhibit lower detectability and are reported as measured operating points, while the security-aware ranker is retained as a future optimization direction. High-payload detectability remains a true limitation for covert-use claims.

\manuscriptchange{The Abstract, Experimental Scope, and Conclusion now include the high-payload detectability limitation, and the security-aware ranker equation is presented as future work.}

\begin{reviewercomment}
4. Figures 5 and 6 should include statistical variance such as 95\% confidence intervals.
\end{reviewercomment}

\response We agree. The multi-payload and protocol-overview figures were regenerated with 95\% confidence information for the payload and distortion views.

\manuscriptchange{The multi-payload and protocol-overview figures now include uncertainty information in the plotted comparisons.}

\section*{Response to Reviewer \#6}

\begin{reviewercomment}
1. The abstract is overloaded with performance numbers and obscures the core innovation.
\end{reviewercomment}

\response We agree. The abstract now states the core mechanism first and reports only the numbers needed to substantiate the central activation effect, frozen external transfer, and operating envelope.

\manuscriptchange{The Abstract was rewritten around serialization-aware activation, exact wire accounting, and the bounded interpretation of the FUNSD evidence.}

\begin{reviewercomment}
2. Related work is too long and tangential.
\end{reviewercomment}

\response We agree. The related work section has been shortened and reorganized around binary pattern/block substitution, distortion and coding, neighbouring domains, and the resulting research gap.

\manuscriptchange{The Related Work section was rewritten and narrowed.}

\begin{reviewercomment}
3. Add more recent state-of-the-art comparisons.
\end{reviewercomment}

\response We expanded the literature positioning with additional relevant binary-cover and side-information-aware works. For the executable common-protocol comparison, we retained ABM, PPOCP, Dong, and Huynh--Nguyen because these are the methods for which a meaningful plaintext binary-image block/pattern RDH comparison could be implemented under matched messages and recovery accounting.

\manuscriptchange{The comparative related-work table and the bibliography now include additional recent and conceptual prior art, while the common protocol explains why only executable direct comparators are included.}

\begin{reviewercomment}
4. The two-layer CNN lacks justification and comparison with deeper architectures.
\end{reviewercomment}

\response We agree. The new architecture diagnostic compares shallow, current, deeper, and window-size variants.

\manuscriptchange{The CNN-architecture table adds this comparison.}

\begin{reviewercomment}
5. The lossless-channel assumption ignores JPEG compression and noise.
\end{reviewercomment}

\response We agree. We added a channel stress test. PNG and TIFF preserve message and cover exactly. JPEG recompression and one-pixel noise break exact extraction/restoration. The revised manuscript now classifies this correctly: preservation of the matched stego/auxiliary pair over a lossless path is a condition of validity of the exact-recovery contract. Robust operation under lossy/noisy/geometrically transformed channels would require an external synchronization or error-control layer and is not claimed here.

\manuscriptchange{The channel-stress table and the ``Conditions of validity'' paragraph in Experimental Scope state the channel requirement explicitly.}

\begin{reviewercomment}
6. The keywords are too generic.
\end{reviewercomment}

\response We agree. The keywords were replaced with more specific terms.

\manuscriptchange{The keyword list now includes binary-image reversible data hiding, PEAK/ZERO block mapping, serialization-aware mapping selection, auxiliary-information coding, and learned distortion ranking.}

\section*{Response to Reviewer \#7}

\begin{reviewercomment}
1. The manuscript requires careful proofreading.
\end{reviewercomment}

\response We agree. We proofread the revised manuscript while rewriting the abstract, introduction, related work, method, experiments, discussion, and conclusion.

\manuscriptchange{Language and flow were edited throughout the manuscript.}

\begin{reviewercomment}
2. Equations are not numbered.
\end{reviewercomment}

\response We agree. The displayed equations used for the method and discussion were converted to numbered equation environments.

\manuscriptchange{The revised source contains numbered equations for the codec contract, block indexing, adaptive threshold, candidate set, net capacity, enumerative state, complexity, and security-aware future objective.}

\begin{reviewercomment}
3. Study the influence of parameters.
\end{reviewercomment}

\response We agree. Payload, auxiliary length, serialization credit, mapping activation sensitivity, and binarization sensitivity are now reported explicitly. The FUNSD fixed-128/Otsu validation further shows how external document preprocessing affects availability, auxiliary cost, net payload, and DRD without changing the exact-recovery result on available runs.

\manuscriptchange{The revised experiments add the activation ablation, serialization sensitivity, fixed-cohort analysis, FUNSD fixed-128/Otsu validation, and payload diagnostic discussion.}

\begin{reviewercomment}
4. Clarify the improvement of the proposed method in the experimental evaluation.
\end{reviewercomment}

\response We agree. The central document-image ablation shows the isolated effect: compared with the same all-feasible mapping family charged after selection, serialization-aware activation reduces mean auxiliary cost and raises mean net payload from 2139 to 2245 bits while preserving exact reversibility. The independent FUNSD validation confirms the direction of this effect: A2 is never worse than A1 in paired net payload and saves approximately 239--241 auxiliary bits per paired image under both tested binarizations.

\manuscriptchange{The activation-ablation table, FUNSD external-validation table, and surrounding text isolate and externally validate the improvement due to the proposed selection rule.}

\begin{reviewercomment}
5. Neural networks are widely used; refer to more state-of-the-art works and add discussion.
\end{reviewercomment}

\response We expanded the neural-network discussion where it is relevant to RDH ranking and steganalysis. Neural-network references from different application domains fall outside the revised citation scope because they address separate tasks from binary-image RDH, reversible coding, side-information cost, and distortion ranking.

\manuscriptchange{The Related Work section now includes learned RDH and steganalysis references relevant to this manuscript's learning component.}

\section*{Response to Reviewer \#8}

\begin{reviewercomment}
1. Describe CNN design, 9-by-9 window rationale, loss function, partitioning, and LOSO validation in more detail.
\end{reviewercomment}

\response We agree. The revised method and experimental sections explain that the 9-by-9 support covers the DRD neighbourhood around a $3\times3$ replacement, describe the regression target and loss, and state the leave-one-source-out image split. The new FUNSD experiment additionally evaluates the committed model under a true frozen external protocol with no training or tuning on the external corpus.

\manuscriptchange{Sections ``Learned Ranking, Ablations, and Wire Accounting,'' ``Machine-Learning Validation,'' and ``Frozen External Validation on FUNSD'' now include these details.}

\begin{reviewercomment}
2. Distinguish contributions from Huynh--Nguyen and Dong, and provide component ablations.
\end{reviewercomment}

\response We agree. The manuscript now distinguishes direct prior art from the revised contribution. Huynh--Nguyen is identified as the closest direct $3\times3$ block-mapping comparator, and Dong is identified as prior art for side information and compression. The central distinction is the use of exact serialized mapping cost in the active mapping-selection decision. The component ablation verifies that this decision changes net payload independently of the ranker.

\manuscriptchange{The Related Work section, ``Comparison with Conventional Block-Mapping'' subsection, and activation-ablation table now make this distinction explicit.}

\begin{reviewercomment}
3. Add targeted detectability analysis explaining why AUC increases.
\end{reviewercomment}

\response We agree. The revised manuscript reports histogram-displacement and transition-change diagnostics. These show that detectability rises with payload because block-histogram mass is shifted from PEAK patterns to ZERO patterns and local transitions increase with changed-pixel count.

\manuscriptchange{Section ``Grouped Steganalysis'' and Section ``Security-Aware Optimization Direction'' now provide this quantitative explanation.}

\begin{reviewercomment}
4. Validation uses only eight document images and 100 thresholded BOSS images; more representative real-world binary images are not considered.
\end{reviewercomment}

\response We agree, and this concern is now directly addressed rather than only acknowledged. We added a frozen external validation on all 199 real scanned FUNSD documents. The committed CNN ranker is applied without retraining, fine-tuning, or tuning on FUNSD, under both fixed-128 and Otsu binarization and at four payload targets. Every available embedding recovers the message and cover exactly. Availability is 98.99--100\%, and the serialization-aware A2 activation is never worse than A1 in paired net payload while saving approximately 239--241 auxiliary bits on average. These results materially weaken the previous dataset-generalization limitation by adding a sizeable independent real scanned-document corpus.

We retain a narrower limitation: FUNSD is one external scanned-form corpus, so transfer across additional independent document collections, languages, scanner/acquisition pipelines, and document genres remains unverified. We also report that mean net payload remains negative through 512 bits on FUNSD; this is an observed operating characteristic of the tested payload/serialization regime, not a failure of exact recovery.

\manuscriptchange{The Datasets subsection, new ``Frozen External Validation on FUNSD'' subsection, Experimental Scope, and Conclusion now replace the former ``larger real-world document corpus'' limitation with measured FUNSD evidence and a narrower remaining external-breadth limitation.}

\section*{Closing Statement}

We believe the revised manuscript now answers the central concerns raised across the reviews: it uses a focused serialization-aware activation claim; defines the exact auxiliary-stream contract; proves the core incremental value through a direct activation ablation; separates deployed components from diagnostics; adds simpler ranker and CNN architecture baselines; adds frozen external validation on 199 real scanned FUNSD documents without retraining; narrows the security claim; and distinguishes true limitations from conditions of validity and operating characteristics. The remaining limitations are stated explicitly without treating metadata-dominated low-payload behavior or the lossless matched-pair contract as generic methodological failures.

\vspace{1.5cm}
Sincerely,

St\'ephane Ga\"el R. Ekodeck and co-authors

\end{document}